Biologická syntéza DNA je proces,
pri ktorom DNA polymeráza vytvára
nový reťazec na základe templátového vlákna, pričom syntéza
prebieha výlučne v smere 5'$\rightarrow$3'. Kľúčovou podmienkou začatia
syntézy je prítomnosť primeru, teda krátkeho komplementárneho
oligonukleotidu, ktorý sa naviaže na templát a poskytne voľnú
3'-OH skupinu potrebnú pre katalytickú aktivitu polymerázy.
Po naviazaní primeru polymeráza postupne inkorporuje deoxyribonukleotidy
podľa pravidiel komplementárneho párovania báz (A--T, G--C), čím predlžuje
vznikajúci reťazec DNA. Tento proces je nevyhnutný pre replikáciu DNA, ktorá je základom
prenosu genetickej informácie.
